\documentclass[conference]{IEEEtran}
\usepackage[utf8]{inputenc}
\usepackage[T1]{fontenc}
\usepackage{cite}
\usepackage{amsmath,amssymb}
\usepackage{graphicx}
\usepackage{booktabs}
\usepackage{url}

\title{A Preregistered Audit of Dataset Contamination Controls in LLM-for-NetOps Benchmarks}

\author{
\IEEEauthorblockN{Autonomous AI Researcher}
\IEEEauthorblockA{CNSM 2026 Agentic AI Researcher Track}
}

\begin{document}

\maketitle

\begin{abstract}
We preregistered and executed a paired confirmatory audit (study\_id = contamination-audit-netops-2026-v1) to evaluate whether a conservative detector suite (exact normalized token equality, fuzzy 5-gram overlap, provenance heuristics) would flag items in a reproducible 150-pair NetOps sample and whether applying a preregistered guarded exclusion policy would change validator-based success rates. The preregistered design fixed detectors and thresholds (fuzzy 5-gram high-confidence = 0.80), sampling (150 pairs), and the primary estimand (paired\_success\_rate\_difference\_guarded\_minus\_baseline). The run used provenance-stable synthetic tasks and shared\_initial\_candidate semantics (independent\_condition\_generation = false), producing one shared generation per item reused across baseline and guarded episodes. No preregistered high-confidence contamination flags occurred; therefore no guarded exclusions or replacements were applied. Validator outcomes were identical (baseline\_success = 150/150; guarded\_success = 150/150), yielding a paired difference of 0.0 percentage points, paired bootstrap 95\% percentile CI [0.0, 0.0] (10,000 resamples, bootstrap\_seed = 1729), and exact McNemar two-sided p = 1.0. We classify the preregistered primary estimand as degenerate/untestable for generation-level contamination under the executed semantics and provide an artifact-grounded reconciliation and concrete design changes (public provenance sampling, independent per-condition generation, full per-item overlap persistence, multi-model replication, paraphrase-robust detectors) to increase audit sensitivity. All execution and analysis artifacts are archived in the execution bundle and summarized in the Disclosure Statement.
\end{abstract}

\section{Introduction}
LLMs are increasingly proposed for network-operations (NetOps) tasks such as intent translation, automated repair, and operational automation. For operational trust, measured evaluation performance must reflect genuine model capability rather than dataset contamination (benchmark items present in model training data) or trivial templating artifacts. Methodological literature therefore recommends layered contamination defenses---deterministic fingerprinting, fuzzy n-gram overlap, provenance checks---and coupling generation with deterministic validators to provide functional oracles. We preregistered and executed a narrowly scoped confirmatory audit to answer: under a conservative, preregistered detector suite and a deterministic guarded exclusion policy, how often would items in a reproducible 150-pair NetOps sample be flagged as high-confidence contamination, and how much would applying those controls change a validator-based success rate for a hosted LLM (declared model: gpt-5-mini)? The preregistration fixed detectors, thresholds (fuzzy 5-gram high-confidence = 0.80), sample size (150 pairs), and the primary paired estimand. This paper reports the artifact-verified execution, exact confirmatory statistics, and an execution--to--preregistration reconciliation that identifies why the primary paired test was uninformative under the executed semantics and prescribes concrete, artifact-anchored redesigns to increase audit sensitivity.

\section{Related Work}
Benchmarks and systems have demonstrated that LLMs can generate network configuration artifacts from high-level intent and that verification loops (generate--validate--repair) reduce regressions compared with one-shot generation [1,2]. Deterministic validators based on header-space analysis and related formal methods provide practical oracles for reachability and policy checks and are appropriate for functional scoring in NetOps workflows [3,4]. Methodological work on evaluation validity highlights risks from dataset contamination and templating and recommends layered detection and conservative exclusion policies; surveys and frameworks call for evidence-bounded, contamination-aware benchmark practices and richer reporting [5,6,7]. Our audit implements a conservative subset of these defenses (exact normalized equality; fuzzy 5-gram overlap with preregistered thresholds 0.50/0.65/0.80; provenance heuristics) and pairs them with a deterministic NetOps validator to ask whether these preregistered controls materially alter measured success on a frozen confirmatory sample.

\section{Methodology}
Preregistration and primary estimand. The preregistration (study\_id contamination-audit-netops-2026-v1) specified a paired design comparing baseline vs guarded conditions on the same item\ensuremath{\times}model pairs with primary estimand paired\_success\_rate\_difference\_guarded\_minus\_baseline. Confirmatory sample size was 150 pairs; detectors were exact normalized token equality, fuzzy 5-gram shingle overlap (high-confidence threshold >= 0.80; sensitivity brackets at 0.65 and 0.50), and provenance timestamp heuristics. The guarded policy deterministically excluded high-confidence flagged items and replaced them with deterministic retemplated augmentations subject to an embedding-similarity safeguard (minimal cosine >= 0.85). Analysis planned an exact two-sided McNemar test and a paired bootstrap percentile 95\% CI with 10,000 resamples (bootstrap\_seed = 1729). All choices and thresholds were fixed prior to execution.

Execution semantics and sampling. The execution used a deterministic task generator to produce provenance-stable synthetic confirmatory items (each labeled with synthetic source identifiers) and applied shared\_initial\_candidate semantics: independent\_condition\_generation = false, meaning a single generation per item could be reused across baseline and guarded episodes unless a high-confidence contamination flag triggered replacement. This reduced hosted-model calls to model\_calls\_used = 150 (the preregistered maximum was 300). The NetOps verification oracle was a deterministic validator returning a binary pass/fail with full deterministic traces (parsed\_command\_count, normalized\_configuration, final\_state, violation codes).

Detector implementation and persistence. The detector suite implemented deterministic heuristics: exact normalized token equality; fuzzy 5-gram shingle overlap with preregistered high-confidence flag at overlap >= 0.80; provenance timestamp comparisons against archive cutoffs. To conserve storage, the executed workflow persisted detailed per-item fuzzy-overlap outputs only when threshold crossings or replacements occurred. The guarded policy excluded only high-confidence flagged items and would deterministically replace them with retemplated items that pass the embedding-similarity safeguard and functional checks; in the executed run no such exclusions occurred and therefore no replacements were generated or persisted.

Analysis executor and reproducibility controls. Confirmatory analysis used the registered executor paired\_binary\_analysis\_v1 producing paired contingency counts, an exact McNemar two-sided p-value, and a paired bootstrap percentile 95\% CI (10,000 resamples, seed = 1729). The execution manifest, raw results, deterministic reconciliation, and analysis artifacts (paired\_contingency\_table.csv, condition\_summary.csv, contamination\_summary.csv) are archived in the execution bundle; the analysis reported below quotes values verbatim from those artifacts.

\section{Results}
Execution accounting and detector outcomes. Execution completed the planned episodes (planned = 300 episodes corresponding to 150 item pairs) and recorded model\_calls\_used = 150 consistent with shared\_initial\_candidate semantics. The detector suite produced zero preregistered high-confidence contamination flags in the executed sample; contamination\_summary.csv is empty and no guarded replacements were applied.

Validator outcomes and confirmatory statistics. The deterministic NetOps validator returned pass for all baseline episodes (150/150) and for all guarded episodes (150/150). The paired contingency counts from analysis/paired\_contingency\_table.csv are n11 = 150, n10 = 0, n01 = 0, n00 = 0. The primary estimand (paired\_success\_rate\_difference\_guarded\_minus\_baseline) = 0.0 percentage points. Paired bootstrap percentile 95\% CI (10,000 resamples, bootstrap\_seed = 1729) = [0.0, 0.0]. Exact McNemar two-sided p = 1.0. These values are reported verbatim from analysis/results.json.

Representative diagnostics. Representative validator traces show shared\_initial\_candidate reuse across conditions when no exclusion applied. For the example items archived in execution/scoring, parsed\_command\_count values observed in representative traces include 2, 4, and 6 (illustrative tasks in the archive), and the validator reported VALID\_CONFIGURATION with explicit normalized\_configuration strings identical between baseline and guarded episodes. Deterministic reconciliation artifacts (analysis/deterministic\_reconciliation.json) confirm completed\_episode\_count = 300, completed\_hosted\_model\_call\_event\_count = 150, and all internal consistency checks passed.

\section{Discussion}
Interpretation of the confirmatory null. The observed p-value and bootstrap CI correctly reflect complete concordance between baseline and guarded outcomes in this run. However, under the executed semantics this concordance is mechanically expected: independent\_condition\_generation was false and the pipeline reused identical model outputs across baseline and guarded conditions unless a high-confidence flag triggered replacement. Because the detectors produced zero high-confidence flags, the guarded intervention had no opportunity to alter generations, rendering the primary paired estimand degenerate for generation-level contamination detection in this execution.

Artifact-grounded diagnosis. Two execution facts in the archived manifest explain the degenerate null: (1) the confirmatory sample used provenance-stable synthetic tasks (task\_manifest entries show source\_identifier fields like synthetic-netops:task-000XXX), and (2) independent\_condition\_generation=false (shared\_initial\_candidate semantics), producing a single shared generation per item reused across both conditions. In addition, the executed persistence policy saved fuzzy-overlap outputs only when threshold crossings or replacements occurred; because none occurred, the archive lacks a full per-item overlap table and therefore precludes reconstructing a sensitivity curve from the executed artifacts.

Design sensitivity and concrete recommendations. Detecting modest paired differences requires semantics and instrumentation that permit observable generation-level discordance and richer detectors. Based on the deterministic reconciliation and archived artifacts we recommend these concrete changes for future confirmatory audits: (a) sample provenance-rich public items (the preregistration originally targeted public OpenAlex benchmark sources) rather than synthetic-only items to expose potential pretraining-era overlaps; (b) set independent\_condition\_generation = true so baseline and guarded outputs are independently realized and can diverge absent exclusions; (c) persist full per-item fuzzy-overlap scores for post-hoc sensitivity analysis rather than only threshold-crossing subsets; (d) expand to multi-model replication and multiple generation seeds to capture stochastic leakage and model heterogeneity; and (e) integrate paraphrase-sensitive detectors (semantic-similarity retrieval or paraphrase-resilient scoring) and human adjudication for borderline flags. Each recommendation follows directly from archived facts (shared generation, empty contamination\_summary.csv, absent per-item overlaps) and addresses the mechanical limitations that rendered the primary paired test uninformative in this run.

Statistical and operational implications. The exact McNemar p = 1.0 and bootstrap CI [0.0,0.0] are valid summaries of the observed paired table but do not quantify contamination prevalence in web-sourced corpora or sensitivity to paraphrased memorization. Practically, deterministic validators are effective functional oracles for scoring (they provided unambiguous pass/fail traces here), but auditors seeking to measure contamination-induced performance differences must ensure the experimental semantics enable independent realizations or apply cross-condition randomness and multi-model designs. The present artifacts support only the narrow inference that, under provenance-stable synthetic sampling and shared-generation semantics, the preregistered conservative detectors produced no high-confidence flags and the guarded policy had no measurable effect.

\section{Limitations}
\begin{itemize}
\item Single-model execution (gpt-5-mini) --- results do not generalize across other LLM families, sizes, or training corpora.
\item The executed sample used provenance-stable synthetic/deterministically generated tasks; absence of preregistered high-confidence flags here may not reflect contamination prevalence in real-world, web-crawled benchmarks (the preregistration had targeted public OpenAlex datasets).
\item Generation semantics used shared\_initial\_candidate across baseline and guarded conditions (independent\_condition\_generation = false), producing identical model outputs when no exclusions occurred and thus limiting sensitivity to interventions that rely on independent re-generation.
\item Contamination detection relied on preregistered conservative heuristics (exact match and fuzzy 5-gram overlap thresholds); subtler contamination patterns (paraphrased memorization, partial leakage) may escape these heuristics and require richer provenance, paraphrase-robust detectors, or human adjudication.
\item Passing the deterministic NetOps validator indicates satisfaction of the checked functional constraints but does not imply comprehensive production readiness or absence of semantic regressions outside the validator's scope.
\end{itemize}

\section{Conclusion}
We executed a fully preregistered, artifact-anchored paired audit of conservative contamination detectors for LLM-driven NetOps evaluation. In the reproducible 150-pair synthetic sample we observed zero preregistered high-confidence flags and identical validator pass rates (150/150 baseline, 150/150 guarded). The degenerate null is explained by execution facts---provenance-stable synthetic sampling and shared\_initial\_candidate semantics---and therefore does not imply absence of contamination in public, web-crawled benchmarks. We document deterministic reconciliation checks that passed, report exact confirmatory statistics verbatim from archived analysis artifacts, and recommend concrete, artifact-grounded design changes (independent per-condition generation, full per-item overlap persistence, multi-model replication, paraphrase-robust detectors, provenance-rich public sampling) to increase audit sensitivity in future work. These recommendations and all reported numerical values follow directly from the archived execution and analysis artifacts.

\section*{Disclosure Statement}
This study was executed under the autonomous-run preregistration contamination-audit-netops-2026-v1. Declared LLM: gpt-5-mini. Orchestration adapter and deterministic components included a hosted orchestration adapter, a deterministic task generator, and a deterministic NetOps validator. All generation prompts, model outputs, validator traces, sampling seeds, replacement rules, execution logs, and analysis artifacts produced by the autonomous pipeline are archived in the execution bundle. The autonomous system performed literature discovery, gap selection, preregistration, experiment execution, analysis, manuscript drafting, simulated peer review, and revision after the autonomy lock; no manual human editing of technical content, analyses, references, figures, tables, captions, or this Disclosure Statement occurred after the autonomy lock. Exact immutable initial master-prompt reference (archived): provenance/master\_prompt.txt (SHA-256 = 1872df1e\allowbreak{}1805d2d9\allowbreak{}6940456c\allowbreak{}a016bd66\allowbreak{}5d1d5196\allowbreak{}add77f5a\allowbreak{}cdf1582b\allowbreak{}b39b15ba).

\begin{thebibliography}{99}
\bibitem{ref1} https://arxiv.org/abs/2604.22513
\bibitem{ref2} https://doi.org/10.48550/arxiv.2606.06212
\bibitem{ref3} https://doi.org/10.48550/arxiv.2605.22092
\bibitem{ref4} http://citeseerx.ist.psu.edu/viewdoc/summary?doi=10.1.1.300.8659
\bibitem{ref5} https://doi.org/10.1007/s11704-026-60308-3
\bibitem{ref6} Buğra Alperen Uluırmak, Rıfat Kurban, "EvalSafetyGap: A Hybrid Survey and Conceptual Framework for LLM Evaluation-Safety Failures", 2026
\bibitem{ref7} Nguyen Van Tu, Sukhyun Nam, James Won‐Ki Hong, "Intent‐Based Network Configuration Using Large Language Models", 2024, 10.1002/nem.2313
\bibitem{ref8} Lei Huang, Weijiang Yu, Weitao Ma, Weihong Zhong, Zhangyin Feng, Haotian Wang, Qianglong Chen, Weihua Peng, Xiaocheng Feng, Bing Qin, Ting Liu, "A Survey on Hallucination in Large Language Models: Principles, Taxonomy, Challenges, and Open Questions", 2024, 10.1145/3703155
\end{thebibliography}

\end{document}
